\documentclass[../../../main.tex]{subfiles}


\begin{document}

references are \cite{Nayga2019}, \cite{Kitt2012}, \cite{Ramezani_Masir2013}, \cite{Castro_Neto2009} and \cite{Sun2023}

\subsubsection{Basics of Strain}

For the purpose of introducing strain into the model, I first establish its effect on the coordinates and exchange couplings 

\begin{align}
	\vec{R}_i &\mapsto \vec{R}_i^\prime = \vec{R}_i + \vec{U}(\vec{R}_i) \\
	\vec{\delta}_{ij} &\mapsto \vec{\delta}_{ij}^\prime = \vec{\delta}_{ij} + \vec{U}(\vec{R}_j) - \vec{U}(\vec{R}_i)
\end{align}

\textcolor{blue}{$\vec{\delta}_{ij} = \vec{R}_j-\vec{R}_i$ ;)} \\

$\vec{U}(\vec{R})$ is the displacement field. We approximate that $\vec{U}$ varies slowly with $\vec{R}$:

\begin{gather}
	\vec{U}(\vec{R}_j) = \vec{U}(\vec{R}_i + \vec{\delta}_{ij}) \approx \vec{U}(\vec{R}_i) + \left(\vec{\delta}_{ij}\vec{\nabla}_{\vec{R}}\right) \vec{U}(\vec{R}) \big\rvert_{\vec{R}_i} \\
\end{gather}


It is convenient to redefine 

\begin{equation}
	\left( \vec{\delta}_{ij} \vec{\nabla}_{\vec{R}} \right) \vec{U}(\vec{R})\big\rvert_{\vec{R}_i} = \underline{u} \cdot \vec{\delta}_{ij},
\end{equation}

where $u_{\alpha\beta} = \partial_{\beta}U_{\alpha}$ is the strain tensor. \\

\begin{equation}
	\Rightarrow \vec{\delta}_{ij}^\prime \approx \vec{\delta}_{ij} + \underline{u}\cdot\vec{\delta}_{ij} 
\end{equation}

\textcolor{red}{Here $\underline{u}$ is taken at $\vec{R}_i$, however one can derive the same expression but with $\underline{u}$ taken at $\vec{R}_j$. In what follows I choose the expansion around $\vec{R}_i$ as shown in the main text. This however is arbitrary. Something dependening on $\vec{R}_i$ and $\vec{R}_j$ (or $(\vec{R}_i + \vec{R}_j)/2$) would seem more sensible but would complicate later expressions. I was not able to resolve this yet.} \\

I do the same approximation for the exchange couplings 

\begin{align}
	J(\vec{\delta}_{ij}) \mapsto J(\vec{\delta}_{ij}^\prime) &\approx J\left(\vec{\delta}_{ij} + \underline{u}\cdot\vec{\delta}_{ij}\right) \\
	&\approx J(\vec{\delta}_{ij}) + \left(\underline{u}\cdot\vec{\delta}_{ij}\right) \cdot \left( \vec{\nabla}_{\vec{\delta}}J(\vec{\delta})\big\rvert_{\delta_{ij}}\right)
\end{align}

For the function $J(\vec{\delta})$ I uses the form

\begin{equation}
	J(\vec{\delta}) = J\exp\left[-\beta\left(|\vec{\delta}|/a-1\right)\right] \approx J\left(1 -\beta\left(|\vec{\delta}|/a-1\right)\right),
\end{equation}

, see e.g. \cite{Pereira2009}, where $a$ is the lattice constant and $\beta = -d\ln(J)/d\ln(|\vec{\delta}|)\big\rvert_{a}$ the magneto-elastic coupling. \\

\begin{align}
	\Rightarrow \vec{\nabla}_{\vec{\delta}}J(\vec{\delta})\big\rvert_{\delta_{ij}} \overset{|\vec{\delta}_{ij}| = a}{=} -\frac{J\beta}{a^2} \vec{\delta}_{ij} \notag  \\
	\Rightarrow J(\vec{\delta}_{ij}) \approx J(\vec{\delta}_{ij}) - \frac{J\beta}{a^2}\left( \vec{\delta}_{ij} \cdot \underline{u} \cdot \vec{\delta}_{ij} \right)
\end{align}



Looking at the matrix elments (\ref{mael11} and below), one also requires the expansion of $\exp[i\vec{\delta}_{ij}\vec{K}]$. To summarize:

\begin{align}
	\vec{\delta}_{ij}^\prime &\approx \vec{\delta}_{ij} + \underline{u} \cdot \vec{\delta}_{ij} \\
	J(\vec{\delta}_{ij}^\prime) &\approx J - \frac{J\beta}{a^2}\left( \vec{\delta}_{ij} \cdot \underline{u} \cdot \vec{\delta}_{ij}\right) \\
	e^{i\vec{K} \cdot \vec{\delta}_{ij}^\prime} &\approx e^{i\vec{K}\cdot\vec{\delta}_{ij}} + ie^{i\vec{K}\cdot\vec{\delta}_{ij}}\left( \vec{K} \cdot \underline{u} \cdot \vec{\delta}_{ij}
	\right)
\end{align}


\textcolor{red}{In this approximation the Mean field does not react to the strain at all??} 

\subsubsection{Straining the 0-$\pi$-State}

The prototypical form of the matrix elements is (expanding around $\vec{K}$)

\begin{equation}
	J(\vec{\delta}_{ij})a_{i\sigma}^\dagger a_{i+\delta_{ij};\sigma} \sim A_{K\sigma}^\dagger \left(J(\vec{\delta}_{ij})e^{i\vec{K}\cdot\vec{\delta}_{ij}}(1+i\vec{\delta}_{ij}\cdot\vec{p})\right) A_{K\sigma}
\end{equation}

Expanding up to first order in strain I acquire the following terms 


\begin{align*}
	&Je^{i\vec{K}\cdot\vec{\delta}_{ij}}(1+i\vec{\delta}_{ij}\cdot\vec{p}) & \text{unstrained Hamiltonian/ zeroth order} \\
	& iJe^{i\vec{K}\cdot\vec{\delta}_{ij}} \left( \vec{K} \cdot \underline{u} \cdot \vec{\delta}_{ij} \right) & \text{contribution to $\vec{A}$ of $\mathcal{O}\left( \underline{u}^1\right)$} \\
	& -\frac{J\beta}{a^2} e^{i\vec{K}\cdot\vec{\delta}_{ij}} \left( \vec{\delta}_{ij} \cdot \underline{u} \cdot \vec{\delta}_{ij} \right) &\text{contribution to $\vec{A}$ of $\mathcal{O}\left( \underline{u}^1\right)$}   \\
	& -Je^{i\vec{K}\cdot\vec{\delta}_{ij}}\left( \vec{K} \cdot \underline{u} \cdot \vec{\delta}_{ij} \right)\left(\vec{\delta}_{ij}\cdot\vec{p}\right) &\text{correction to $v_F$ of $\mathcal{O}\left( \underline{u}^1 \right)$}\\
	& -i\frac{J\beta}{a^2}e^{i\vec{K}\cdot\vec{\delta}_{ij}}\left(\vec{\delta}_{ij} \cdot \underline{u} \cdot \vec{\delta}_{ij}\right)\left(\vec{\delta}_{ij}\cdot\vec{p}\right) &\text{correction to $v_F$ of $\mathcal{O}\left( \underline{u}^1 \right)$}\\
	& iJe^{i\vec{K}\cdot\vec{\delta}_{ij}}\left(\underline{u}\cdot\vec{\delta}_{ij}\right)\cdot\vec{p} &\text{correction to $v_F$ of $\mathcal{O}\left( \underline{u}^1 \right)$}
\end{align*}

(I drop the lattice constant again from now on) I will focus on the terms of order zero + terms of either order 1 in field gradient or order 1 in strain \cite{Nayga2019}. \textcolor{red}{I havent come up with a good argument for this yet. How does one even compares orders between different expansions?} This gives

\begin{align}
	\begin{split}
		J(\vec{\delta}_{ij})e^{i\vec{K}\cdot\vec{\delta}_{ij}}(1+i\vec{\delta}_{ij}\cdot\vec{p}) &\sim  Je^{i\vec{K}\cdot\vec{\delta}_{ij}}\left(1 + i\vec{\delta}_{ij}\cdot\vec{p} + i \left(\vec{K} \cdot \underline{u} \cdot \vec{\delta}_{ij}\right) - \beta\left(\vec{\delta}_{ij} \cdot \underline{u} \cdot \vec{\delta}_{ij}\right)\right)\\ 
		&= Je^{i\vec{K}\cdot\vec{\delta}_{ij}}\left(1 + i \vec{\delta}_{ij}\cdot \left[\vec{p} + \underline{u}^T\left(\vec{K} + i \beta\vec{\delta}_{ij}\right)\right]\right)
	\end{split}
\end{align}

Hence, on this level

\begin{equation}
	\vec{A} = \underline{u}^T\left(\vec{K} + i \beta\vec{\delta}_{ij}\right).
\end{equation}

\textcolor{red}{This is not correct since one in has to add contributions from different local directions in the off-diagonal terms of $H$!}

\begin{equation}
	B = \left(\partial_x A_y - \partial_y A_x\right)
\end{equation}

It turns out however, that $\vec{B}=0$, which is explained in what follows. The first term is simply rotationless (this was shown here \cite{Ramezani_Masir2013}) and can therefore be "gauged away". \\

Regarding the second term; I take a step back to the dispersion relation of the unstrained model

\begin{equation}
	H = \sum_{\vec{q}\sigma} \Psi_{q\sigma}^\dagger \cdot \left( \vec{h}(\vec{q})\cdot \vec{\sigma} \right) \cdot \Psi_{q\sigma} \quad \text{with} \quad \vec{h}(\vec{q}) = 2t
	\begin{pmatrix}
		\sin(\vec{q}\cdot\vec{\delta}_3) \\
		\cos(\vec{q}\cdot\vec{\delta}_2) \\
		\cos(\vec{q}\cdot\vec{\delta}_1)
	\end{pmatrix} \Rightarrow \varepsilon(\vec{q}) = \abs{\vec{h}(\vec{q})}.
\end{equation}

By definition, $\varepsilon(\vec{K}^{(\prime)}) =0$. For this to be true, every component of $\vec{h}$ must vanish individually at $\vec{K}^{(\prime)}$. This gives three conditions for $\vec{q}$, one for each direction. \\

Evidently each expression in $\vec{h}$ is obtained by adding contribtions from hopping forwards and backwards along $\vec{\delta}_{ij}$ . However, since $\vec{\delta}_{ij} \cdot \underline{u} \cdot \vec{\delta}_{ij}$ is invariant under $\vec{\delta}_{ij} \mapsto -\vec{\delta}_{ij}$ it acquires the same dispersion-like prefactor as the zeroth order, which vanishes by definition at $\vec{K}^{(\prime)}$. \\

So why does the term vanish? Its a mix of the special structure of the hamiltonian (i.e. the local directions decouple in the pauli-matrix basis) and the fact that the expansion term $\vec{\delta}_{ij} \cdot \underline{u} \cdot \vec{\delta}_{ij}$ cant distinguish between $\vec{\delta}_{ij}$ and $-\vec{\delta}_{ij}$. The latter statement is obviously not model-specific. \textcolor{blue}{The same mechanism shouild be present in the $\pi$-flux square lattice. For uniform hopping (e.g. ordinary Graphene) it is absent (the decoupling of the directions is a cosnequence of the non-uniform hopping/ the mean-field).} \\

\textcolor{red}{It is not clear to me how this can be upheld under gauge transformations (wouldn't this mix directions in the Pauli-matrix basis?)} \\

\textcolor{blue}{I would expect to find a pseudo-magnetic field in higher orders of strain?}

\end{document}
